% !TeX root = ../../Book1.tex
\chapter{泛函分析补充}
\label{b1:app:C}

\begin{chapterabstract}
本附录围绕有界线性方程的可解性补充泛函分析工具。
商空间去掉解的不唯一方向，对偶算子给出相容条件，闭值域估计把近似可解提升为精确可解。
紧扰动使核与余核降为有限维；在自伴情形，谱分解进一步给出解的坐标与近似误差。
最后将凸分离用于统一间隙、距离公式和最佳逼近。
\end{chapterabstract}

\begin{learninggoals}
\begin{enumerate}
\item\label{b1:item:C1-goal-quotient}用商范数表述稳定性，区分 Banach 空间对偶算子与 Hilbert 空间伴随。
\item\label{b1:item:C1-goal-closed-range}从对偶消去条件得到近似可解性，并说明闭值域估计怎样补足存在性。
\item\label{b1:item:C1-goal-fredholm}证明恒等算子的紧扰动具有有限维核、闭值域和零指标。
\item\label{b1:item:C1-goal-spectral}使用已有紧自伴谱定理求解方程，辨认共振、小特征值及截断稳定性。
\item\label{b1:item:C1-goal-separation}将凸分离转成定量距离估计，辨明最近点存在与对偶最大值取得的区别。
\end{enumerate}
\end{learninggoals}

有限维线性代数中，检查右端是否满足左零空间给出的约束，就能判定方程有没有解。
搬到函数空间以后，这个判断首先只能得到右端在值域的闭包中。
值域可能不闭，近似解也可能没有任何有界子列。分析中的先验估计所承担的任务之一，
便是阻止这一失败；它需要估计全部解中去掉齐次解以后的部分，而不一定控制任意原像。

读者应熟悉\chapref{b1:ch:08}的 Hahn--Banach 定理、开映射定理、Hilbert 投影和紧算子，
以及\chapref{b1:ch:09}的弱拓扑与自然嵌入。商空间完备性及其对偶已经在%
\cref{b1:ex:08-quotient-dual}证明，双零化子关系见%
\cref{b1:ex:09-double-annihilator}。
本附录沿这些结果继续，不重新把相关习题改写为一遍正文。
若主要关心 Hilbert 空间方程，可先掌握商空间估计，再阅读闭值域的正交形式和谱应用；
一般 Banach 空间的 Fredholm 证明则说明哪些结论并不需要自伴性。

Hunter 与 Nachtergaele 的《Applied Analysis》\cite[Section 8.3]{Hunter2001Applied}
适合对照伴随、值域与可解性条件的关系；该书采用的内积线性变量与本书不同，
读公式时应以定义为准。本附录把可解性推导推进到一般 Banach 空间，并给出所用紧扰动结论的证明。
已有许全华、马涛、尹智的《泛函分析讲义》\cite{Xu2017Fanhan}可作为中文衔接读物：
第 9.1、9.2 节的“共轭算子”“商空间”是核对名称和查找相关内容的入口。
若要从这里系统进入一般算子谱与局部凸空间，
可继续阅读 Rudin 的《Functional Analysis》\cite{Rudin1991Functional}；
本附录不涉及一般无界算子的定义域理论，也不以有界算子的公式代替那些额外问题。

% !TeX root = ../../Book1.tex
\section{商空间与对偶算子}
\label{b1:sec:C101}

求解线性方程 \(Tx=y\) 时，两个不同的向量可能给出同一个像。这种不唯一性完全由
\(\ker T\) 描述；若先把相差一个核向量的点视为同一点，剩下的问题便成为单射算子的求逆。
商空间因此不只是代数上的简写。商范数还会告诉我们：在忽略不唯一部分以后，一个近似解究竟离真解多远。

\subsection{商映射与算子的因子分解}
\label{b1:subsec:C10101}

设 \(M\) 为 Banach 空间 \(X\) 的闭线性子空间。沿用%
\cref{b1:thm:quotient-banach}中的商范数与商映射：
\[
 \|x+M\|_{X/M}=\operatorname{dist}(x,M),\quad Qx=x+M.
\]
该定理及\cref{b1:prop:quotient-dual}已经证明 \(X/M\) 完备，并把 \((X/M)^*\) 等距识别为
\[
 M^\circ=\{\,f\in X^*:f(m)=0\ \text{对每个 }m\in M\,\}.
\]
本节直接使用这两个正文结果。代数商空间 \(X/M\) 的定义本身不要求 \(M\) 闭；
不过，要使上式给出的商半范数成为范数，需要 \(M\) 闭。若 \(M\) 不闭，
它会把 \(\overline M\) 中的所有点也压成距离零。本节始终只考虑闭子空间。

\ProseParagraphBreak

商映射不仅连续，而且把开球精确地映为开球。对 \(r>0\)，有
\[
 Q\bigl(B_X(x,r)\bigr)=B_{X/M}(Qx,r).
\]
一个方向来自 \(\|Qz\|\leq\|z\|\)。另一个方向使用严格不等式：
若 \(\|v-Qx\|<r\)，按下确界定义可为这个商元素选取范数小于 \(r\) 的代表 \(z\)，
于是 \(v=Q(x+z)\)。闭球的像却未必等于闭球，因为距离的下确界未必取得。
下面所有选代表的步骤都会保留一个任意小的余量，而不假定存在最短代表。

\begin{proposition}[商空间上的因子分解]\label{b1:prop:C1-quotient-factor}
设 \(X,Y\) 为 Banach 空间，\(M\subset X\) 为闭线性子空间，
\(T\in\mathcal B(X,Y)\) 且 \(M\subset\ker T\)。则存在唯一的
\(\widetilde T\in\mathcal B(X/M,Y)\)，使 \(T=\widetilde TQ\)，并且
\(\|\widetilde T\|=\|T\|\)。若取 \(M=\ker T\)，则 \(\widetilde T\) 单射，
其值域等于 \(\operatorname{ran}T\)。
\end{proposition}
\begin{proof}
定义 \(\widetilde T(x+M)=Tx\)。更换代表时差属于 \(M\)，被 \(T\) 消去，
故定义良好；线性和唯一性由 \(Q\) 满射得到。对任意 \(m\in M\)，
\[
 \|\widetilde T(x+M)\|=\|T(x+m)\|\leq\|T\|\cdot\|x+m\|.
\]
取下确界给出 \(\|\widetilde T\|\leq\|T\|\)。反过来，
\(\|Tx\|\leq\|\widetilde T\|\cdot\|Qx\|\leq\|\widetilde T\|\cdot\|x\|\)，
故两范数相同。当 \(M=\ker T\) 时，\(\widetilde T(x+M)=0\) 恰好表示 \(x+M=0\)。
\end{proof}

这里的单射结论还没有给出连续逆。若把值域赋予 \(Y\) 的子空间范数，则
\(\widetilde T^{-1}\) 有界恰好意味着存在 \(C\)，使
\[
 \operatorname{dist}(x,\ker T)\leq C\|Tx\|\quad(x\in X).
\]
这项估计允许核非零；它估计的是解的等价类，而不是任意选出的一个解。
对于微分方程，核中的向量常是常数或有限个齐次解，因此这种写法比要求
\(\|x\|\leq C\|Tx\|\) 更贴近问题本身。

\subsection{对偶算子的方向与标量}
\label{b1:subsec:C10102}

以下 \(X^*\) 始终由连续线性泛函组成，标量域为 \(\mathbb K=\mathbb R\) 或 \(\mathbb C\)。
若 \(g\) 测试输出 \(Tx\)，那么 \(g\circ T\) 就是对输入的测试。这个复合运算把
算子的方向反转，却不需要在 \(X,Y\) 上引入内积。

\begin{definition}\label{b1:def:C1-dual-operator}
设 \(T\in\mathcal B(X,Y)\)。规定
\[
 T':Y^*\longrightarrow X^*,\quad (T'g)(x)\coloneqq g(Tx),
\]
并称 \(T'\) 为 \(T\) 的\term[duiou-suanzi]{对偶算子}[dual operator]。%
\footnote{许全华等《泛函分析讲义》\cite{Xu2017Fanhan}第 9.1 节使用
“共轭算子”这一名称。\chapref{b1:ch:08}已经说明同一中文名称也用于 Hilbert 空间伴随；
本附录用 \(T'\) 指作用在连续对偶上的算子，用 \(T^*\) 保留 Hilbert 空间伴随，
阅读其他教材时须先看定义域与值域。}
\termidx[gonge-suanzi]{共轭算子}
\end{definition}
\symidx{\(T'\)：有界线性算子的对偶算子}

\begin{proposition}\label{b1:prop:C1-dual-calculus}
对偶算子是有界线性算子，且 \(\|T'\|=\|T\|\)。若复合有意义，则
\[
 (ST)'=T'S',\quad (aT+bU)'=aT'+bU',\quad I'=I.
\]
自然嵌入 \(J_X:X\rightarrow X^{**}\) 满足 \(T''J_X=J_YT\)。
算子 \(T'\) 对相应的\WeakStar{}拓扑连续。
\end{proposition}
\begin{proof}
由
\[
 |(T'g)(x)|=|g(Tx)|\leq\|g\|\cdot\|T\|\cdot\|x\|
\]
得 \(\|T'\|\leq\|T\|\)。对每个 \(x\)，由%
\cref{b1:cor:norming-functional}，
\[
 \|Tx\|=\sup_{\|g\|\leq1}|g(Tx)|
       =\sup_{\|g\|\leq1}|(T'g)(x)|
       \leq\|T'\|\cdot\|x\|,
\]
从而反向不等式成立。复合与线性组合等式逐个在 \(x\) 上求值即可得到。
同样，对 \(g\in Y^*\)，有
\[
 (T''J_Xx)(g)=(J_Xx)(T'g)=g(Tx)=(J_YTx)(g).
\]
最后，若 \(g_\alpha\) \WeakStar{}收敛于 \(g\)，则对每个固定的 \(x\)，
\((T'g_\alpha)(x)=g_\alpha(Tx)\rightarrow g(Tx)\)。这正是像网\WeakStar{}收敛的定义。
\end{proof}

在复 Hilbert 空间中，令 \(R_Hh=\langle\,\cdot,h\rangle\)。按本书第一变量线性的约定，
\(R_H:H\rightarrow H^*\) 是共轭线性等距映射，而
\[
 T'R_Y=R_XT^*.
\]
因此 \((aT)'=aT'\)，但 \((aT)^*=\overline a\,T^*\)。
两条公式没有矛盾：从对偶空间回到 Hilbert 空间时，Riesz 对应本身就共轭了标量。
例如 \(Tz=iz\) 在一维复空间上的对偶算子仍是乘 \(i\)，Hilbert 伴随却是乘 \(-i\)。
把 \(T'\) 与 \(T^*\) 不加说明地混写，会在复系数方程中改变相容条件。

\subsection{核与值域的对偶关系}
\label{b1:subsec:C10103}

对 \(G\subset X^*\)，记
\[
 {}^\circ G=\{\,x\in X:f(x)=0\ \text{对每个 }f\in G\,\}.
\]
上标圆圈在集合右侧时取连续对偶中的泛函，在左侧时取原空间中的向量。
这一记号只是把“被全部测试消去”的条件集中写出，不把 \(X\) 与 \(X^{**}\) 视为同一空间。

\begin{proposition}\label{b1:prop:C1-range-annihilator}
设 \(T\in\mathcal B(X,Y)\)。则
\[
 \ker T'=(\operatorname{ran}T)^\circ,\quad
 \ker T={}^{\circ}(\operatorname{ran}T'),\quad
 \overline{\operatorname{ran}T}={}^{\circ}(\ker T').
\]
最后一式中的闭包是范数闭包。因此 \(T\) 的值域稠密，当且仅当 \(T'\) 单射。
\end{proposition}
\begin{proof}
第一个等式把 \(T'g=0\) 展开为 \(g(Tx)=0\) 对所有 \(x\) 成立。
若 \(Tx=0\)，则全部 \(T'g\) 都消去 \(x\)；反过来，若 \(g(Tx)=0\) 对所有
\(g\in Y^*\) 成立，由连续泛函分离点得到 \(Tx=0\)。这给出第二个等式。
对第三个等式，先由第一个等式得到闭包包含于右边；若 \(y\) 不在值域闭包，
对闭线性子空间 \(\overline{\operatorname{ran}T}\) 使用%
\cref{b1:prop:double-annihilator}中的分离结论，可取 \(g\) 消去这个子空间而
\(g(y)\ne0\)。此时 \(g\in\ker T'\)，故 \(y\) 不在右边。
\end{proof}

该命题只检测近似可解性。条件 \(g(y)=0\) 对所有齐次对偶解 \(T'g=0\) 成立，
保证 \(y\) 可以由 \(Tx_n\) 逼近，但并未保证原像 \(x_n\) 不逃向无穷远。
要从近似解得到真解，还需要值域的闭性。这个缺口正是后面闭值域定理所处理的内容。

% !TeX root = ../../Book1.tex
\section{紧算子及其对偶}
\label{b1:sec:C102}

\chapref{b1:ch:08}已给出紧算子的定义及其在复合、相加和算子范数极限下的稳定性。
本节补充两个后面需要的事实：无穷维空间的单位球不紧；紧性可以在算子与对偶算子之间传递。
前者用于把紧性转化成有限维性，后者使原方程与对偶方程能接受同样的分析。

\subsection{离开闭子空间的单位向量}
\label{b1:subsec:C10201}

\begin{lemma}[Riesz]\label{b1:lem:C1-riesz}
设 \(M\) 为赋范空间 \(X\) 的真闭线性子空间，\(0<\theta<1\)。
则存在单位向量 \(x\in X\)，使 \(\operatorname{dist}(x,M)>\theta\)。
\end{lemma}
\begin{proof}
取 \(z\notin M\)，并令 \(d=\operatorname{dist}(z,M)>0\)。
选取 \(m_0\in M\)，使 \(\|z-m_0\|<d/\theta\)，再置
\(x=(z-m_0)/\|z-m_0\|\)。对每个 \(m\in M\)，有
\[
 \|x-m\|
 =\frac{\bigl\|z-\bigl(m_0+\|z-m_0\|m\bigr)\bigr\|}{\|z-m_0\|}
 \geq\frac d{\|z-m_0\|}>\theta.
\]
取下确界即得所需估计。
\end{proof}

这就是\term[riesz-yinli]{Riesz 引理}[Riesz lemma]。
若 \(X\) 无穷维，可以依次对前面有限多个向量张成的子空间使用该引理，
构造单位向量列，使任意不同两项的距离大于 \(1/2\)。
有限维子空间闭，因此每一步都合法。该列无 Cauchy 子列，单位球便不可能相对紧。
结合有限维空间中有界列的子列定理，得到：赋范空间的闭单位球紧，当且仅当空间有限维。
这里没有假定空间存在内积，也没有使用正交向量。

这个引理中的 \(\theta<1\) 留出了不一定取到最短距离的余量。在 Hilbert 空间中
可以取与 \(M\) 正交的单位向量而得到距离恰为 \(1\)；一般 Banach 空间不能靠“取正交补”
完成这一步。

\subsection{紧性在对偶中的传递}
\label{b1:subsec:C10202}

\begin{theorem}[Schauder]\label{b1:thm:C1-schauder}
设 \(X,Y\) 为 Banach 空间，\(T\in\mathcal B(X,Y)\)。
则 \(T\) 紧，当且仅当 \(T':Y^*\rightarrow X^*\) 紧。
\end{theorem}
\begin{proof}
先设 \(T\) 紧。令 \(C=\overline{T(B_X(0,1))}\)，这是 \(Y\) 中的紧集。
考虑 \(Y^*\) 的闭单位球中一列泛函 \(g_n\)。对每个正整数 \(k\)，选取 \(C\) 的有限
\(1/k\)-网。所有这些网点组成至多可数的集合 \(\{y_j\}\)。
由于 \(|g_n(y_j)|\leq\|y_j\|\)，逐点抽子列并作对角选取，可使
\(g_{n_\ell}(y_j)\) 对每个 \(j\) 收敛。

给定 \(\varepsilon>0\)，选一个网的半径小于 \(\varepsilon/4\)。
在有限个网点上同时使用 Cauchy 条件，便有：对充分大的 \(\ell,m\)，
\[
 \sup_{y\in C}|g_{n_\ell}(y)-g_{n_m}(y)|<\varepsilon.
\]
事实上，先将 \(y\) 移到某个网点，两个泛函各造成不超过网半径的误差；
网点处的差小于 \(\varepsilon/2\)。于是
\[
 \|T'g_{n_\ell}-T'g_{n_m}\|
 =\sup_{\|x\|\leq1}|g_{n_\ell}(Tx)-g_{n_m}(Tx)|<\varepsilon.
\]
连续对偶 \(X^*\) 完备，故这条像子列收敛。任意有界泛函列可先缩放到单位球，
所以 \(T'\) 紧。

反过来，若 \(T'\) 紧，将已证方向用于 \(T'\)，得到 \(T''\) 紧。
设 \((x_n)\) 是 \(X\) 中的有界列。等距性使 \((J_Xx_n)\) 有界，因而可取子列，
使
\[
 T''J_Xx_{n_k}=J_YTx_{n_k}
\]
在 \(Y^{**}\) 中收敛。等距嵌入 \(J_Y(Y)\) 由于 \(Y\) 完备而闭：
其中 Cauchy 列的原像在 \(Y\) 中也为 Cauchy 列，极限仍属于该像。
所以所得极限属于 \(J_Y(Y)\)，原像列 \(Tx_{n_k}\) 在 \(Y\) 中收敛，证明 \(T\) 紧。
\end{proof}

\term[schauder-jinxing]{Schauder 定理}[Schauder theorem]不要求 \(X,Y\) 自反。
反向证明虽然进入了二次对偶，但只使用自然嵌入的等距性和闭像，不能把这一步改写成
“因为 \(Y=Y^{**}\)”。

\subsection{Hilbert 空间中的有限秩逼近}
\label{b1:subsec:C10203}

\chapref{b1:ch:08}证明过用有限秩算子作算子范数逼近足以推出紧性。
在一般 Banach 空间中，逆向结论需要额外条件；若值域是 Hilbert 空间，则有如下直接构造。

\begin{proposition}\label{b1:prop:C1-hilbert-finite-rank}
设 \(X\) 为 Banach 空间，\(H\) 为 Hilbert 空间，\(T:X\rightarrow H\) 紧。
给定 \(\varepsilon>0\)，存在有限维子空间 \(E\subset H\)，使其正交投影 \(P_E\) 满足
\[
 \|T-P_ET\|<\varepsilon.
\]
\end{proposition}
\begin{proof}
在紧集 \(\overline{T(B_X(0,1))}\) 中取有限个半径小于 \(\varepsilon/2\) 的球覆盖，
球心为 \(y_1,\ldots,y_m\)，令 \(E\) 为这些向量张成的空间。
对 \(\|x\|\leq1\)，可取一个球心使 \(\|Tx-y_j\|<\varepsilon/2\)。
正交投影给出最近点，故
\[
 \|Tx-P_ETx\|=\operatorname{dist}(Tx,E)
 \leq\|Tx-y_j\|<\varepsilon/2.
\]
取单位球上的上确界便得到
\(\|T-P_ET\|\leq\varepsilon/2<\varepsilon\)。
\end{proof}

有限维近似由紧像的有限网确定，而不是由 \(X\) 的某个预先给定的坐标截断确定。
即使 \(H\) 不可分，每次所需的近似空间也只有有限维。又因为
\((P_ET)'=T'P_E'\)，这个构造与 Schauder 定理相容，但它没有证明任意紧算子都自伴，
更没有把有限秩近似误当成特征向量分解。

% !TeX root = ../../Book1.tex
\section{闭值域与可解性}
\label{b1:sec:C103}

对于有限维矩阵，值域总是闭的，所以一致线性方程组的极限仍然一致。
无穷维空间中的区别在于：右端可能收敛，所需原像却越来越大。
例如\chapref{b1:ch:08}的对角算子 \(Te_n=n^{-1}e_n\)，可以把单位向量压得任意小。
本节把排除这种现象的条件写成估计，并说明这个估计怎样传到对偶方程。

\subsection{去掉核以后的下界}
\label{b1:subsec:C10301}

\begin{theorem}[闭值域的估计判别]\label{b1:thm:C1-closed-range-estimate}
设 \(X,Y\) 为 Banach 空间，\(T\in\mathcal B(X,Y)\)。以下条件等价：
\begin{equivalences}
\item\label{b1:item:C1-range-closed}\(\operatorname{ran}T\) 在 \(Y\) 中闭；
\item\label{b1:item:C1-quotient-bound}存在 \(C>0\)，使
\[
 \operatorname{dist}(x,\ker T)\leq C\|Tx\|\quad(x\in X).
\]
\end{equivalences}
在这些条件下，对每个 \(y\in\operatorname{ran}T\)，方程 \(Tx=y\) 的解组成一个商类，
而从 \(\operatorname{ran}T\) 到 \(X/\ker T\) 的解映射有界。
\end{theorem}
\begin{proof}
若值域闭，则它是 Banach 空间。由%
\cref{b1:prop:C1-quotient-factor}，
\(\widetilde T:X/\ker T\rightarrow\operatorname{ran}T\) 是有界双射。
有界逆定理给出 \(\|\widetilde T^{-1}y\|\leq C\|y\|\)，即所述估计。

反过来，设估计成立，\(Tx_n\rightarrow y\)。对任意 \(m,n\)，有
\[
 \|Qx_n-Qx_m\|
 =\operatorname{dist}(x_n-x_m,\ker T)
 \leq C\|Tx_n-Tx_m\|.
\]
所以 \(Qx_n\) 在完备的商空间中收敛于某个 \(u\)。连续性给出
\(\widetilde Tu=y\)，因此 \(y\in\operatorname{ran}T\)。这里没有声称原来的
\(x_n\) 收敛；它们可以在核的方向任意漂移，真正收敛的是商类。
\end{proof}

当 \(\ker T=\{0\}\) 时，估计简化为 \(\|Tx\|\geq c\|x\|\)。
它比单射性强：单射性只禁止某个非零向量被压成零，估计则禁止单位向量列的像趋于零。
当核非零时，单位向量本来就可能在核内；正确的归一化是到核的距离等于 \(1\)，
而不是单纯要求 \(\|x\|=1\)。

闭值域保证有界的商类解映射，但不自动给出一个取值于 \(X\) 的有界线性选解算子。
要线性选取代表，需要核有一个闭补空间。有限维核总有这种补空间：取核的一组基，
将其坐标泛函用 Hahn--Banach 延拓到 \(X\)，所得有限秩投影的核就是所需补空间。
Hilbert 空间中则可以统一使用 \((\ker T)^\perp\)。

\subsection{Banach 闭值域定理}
\label{b1:subsec:C10302}

\begin{theorem}[Banach，闭值域]\label{b1:thm:C1-banach-closed-range}
设 \(X,Y\) 为 Banach 空间，\(T\in\mathcal B(X,Y)\)。
则 \(\operatorname{ran}T\) 范数闭，当且仅当 \(\operatorname{ran}T'\) 范数闭。
当它们闭时，
\[
 \operatorname{ran}T={}^{\circ}(\ker T'),\quad
 \operatorname{ran}T'=(\ker T)^\circ.
\]
\end{theorem}
\begin{proof}
先设 \(\operatorname{ran}T\) 闭，并取%
\cref{b1:thm:C1-closed-range-estimate}中的常数 \(C\)。
若 \(f\in(\ker T)^\circ\)，在值域上定义
\(\varphi(Tx)=f(x)\)。因 \(f\) 消去核，此定义与原像选择无关，而且
\[
 |\varphi(Tx)|=|f(x)|
 \leq\|f\|\operatorname{dist}(x,\ker T)
 \leq C\|f\|\cdot\|Tx\|.
\]
Hahn--Banach 定理把 \(\varphi\) 延拓为 \(g\in Y^*\)，于是 \(T'g=f\)。
反向包含 \(\operatorname{ran}T'\subset(\ker T)^\circ\) 不需闭性就成立。
因此 \(\operatorname{ran}T'=(\ker T)^\circ\)，后者是闭子空间。
值域的另一等式由\cref{b1:prop:C1-range-annihilator}去掉闭包得到。

现在设 \(\operatorname{ran}T'\) 闭。将闭值域的估计判别用于 \(T'\)，可取 \(C>0\) 使
\[
 \operatorname{dist}(g,\ker T')\leq C\|T'g\|\quad(g\in Y^*).
\]
记 \(Z=\overline{\operatorname{ran}T}\)。对 \(z\in Z\)，每个 \(h\in\ker T'\) 都消去 \(z\)，故
\begin{equation}\label{b1:eq:C1-dual-lifting-bound}
 |g(z)|\leq\|z\|\operatorname{dist}(g,\ker T')
 \leq C\|z\|\cdot\|T'g\|.
\end{equation}
这个不等式先给出范数可控的近似原像，再由级数把近似变为精确。

具体地，令 \(D=\overline{T(\overline B_X(0,1))}\)，这是 \(Y\) 中闭、凸且平衡的集合。
我们断言 \(\overline B_Z(0,1)\subset CD\)。
若某个 \(\|z\|\leq1\) 的 \(z\in Z\) 不属于 \(CD\)，则凸分离定理给出
\(g\in Y^*\)，使
\[
 \operatorname{Re}g(z)>\sup_{v\in CD}\operatorname{Re}g(v)=C\|T'g\|.
\]
最后的等式使用复平衡性；实标量情形使用正负对称性。
这与\eqref{b1:eq:C1-dual-lifting-bound}矛盾，断言得证。

取 \(z\in Z\)。由刚得到的包含关系，可以选 \(x_1\) 满足
\(\|x_1\|\leq C\|z\|\) 及 \(\|z-Tx_1\|\leq\|z\|/2\)；
若某次余项已经为零，后面全部取零。
否则对余项再次作同样选择，得到
\[
 \|x_j\|\leq C\,2^{-(j-1)}\|z\|,\quad
 \left\|z-T\sum_{j=1}^n x_j\right\|\leq2^{-n}\|z\|.
\]
由 \(X\) 完备，级数 \(x=\sum_jx_j\) 收敛，且 \(\|x\|\leq2C\|z\|\)。
连续性给出 \(Tx=z\)。因此 \(Z\subset\operatorname{ran}T\)，值域闭。
再用证明的前半部分，得到第二个值域等式。
\end{proof}

把\term*[banach-bizhiyu]{Banach 闭值域定理}[Banach closed range theorem]用于
存在性问题时，应分别核对“测试消去”与“值域闭”。
前者往往来自积分分部，后者则来自先验估计、紧性或某种有限维约化。
仅写出形式上的伴随方程，不足以证明原方程有解。

上述反向证明也解释了完备性的位置。对偶估计首先只保证任意小残差，
然后以几何级数控制校正项，原空间的完备性才允许把它们求和。
若直接从 \(z\in\overline{T(B_X)}\) 选出精确原像，就跳过了这一技术桥。

\subsection{Hilbert 空间中的正交解}
\label{b1:subsec:C10303}

\begin{corollary}\label{b1:cor:C1-hilbert-closed-range}
设 \(T:H\rightarrow G\) 为 Hilbert 空间间的有界线性算子。
若 \(\operatorname{ran}T\) 闭，则
\[
 \operatorname{ran}T=(\ker T^*)^\perp,\quad
 \operatorname{ran}T^*=(\ker T)^\perp.
\]
每个满足 \(y\perp\ker T^*\) 的右端有唯一的解 \(x_0\perp\ker T\)；
它是全部解中范数最小的一个，并满足 \(\|x_0\|\leq C\|y\|\)。
\end{corollary}
\begin{proof}
由 Riesz 对应 \(T'R_G=R_HT^*\)，Banach 闭值域定理中的消去条件变为正交条件。
任一解 \(x\) 可唯一分解为 \(x=x_0+n\)，其中 \(x_0\perp\ker T\)，\(n\in\ker T\)。
投影后仍有 \(Tx_0=y\)，两个这样的解相减同时属于核及其正交补，只能为零。
由勾股等式 \(\|x_0+n\|^2=\|x_0\|^2+\|n\|^2\)，最小范数性成立。
最后，\(\|x_0\|=\operatorname{dist}(x_0,\ker T)\)，所以前述估计适用。
\end{proof}

若右端 \(y\) 不满足正交条件，也可以先将其投影到 \(\operatorname{ran}T\)。
此时最小化 \(\|Tx-y\|\) 等价于解
\[
 T^*Tx=T^*y.
\]
这是因为 \(Tx-y\) 与值域正交当且仅当被 \(T^*\) 消去。
只要值域闭，投影存在，从而得到最小二乘解；若值域不闭，最小残差的下确界可能不取得。
闭值域因此同时联系精确可解性、解的稳定性和最小二乘问题。

% !TeX root = ../../Book1.tex
\section{恒等算子的紧扰动与 Fredholm 理论}
\label{b1:sec:C104}

方程 \(x-Kx=y\) 常在先求出某个简单方程的逆以后出现。
算子 \(K\) 若紧，未知量的无限维部分就受到很强的限制：
不唯一的方向只有有限多个，可解性也只留下有限多个线性条件。
这一节不要求 \(K\) 自伴，甚至不要求底空间是 Hilbert 空间。

Hunter 与 Nachtergaele 的《Applied Analysis》\cite[Section 8.3]{Hunter2001Applied}
用原方程与伴随方程说明 Fredholm 理论所描述的可解性模式。
下面将这一路线放到 Banach 空间的对偶配对中，并直接证明恒等算子的紧扰动所需的结论。
读者要区分三个论断：核有限维、值域闭、核与余核维数相同。前两个尚不自动推出第三个。

\subsection{有限维核与闭值域}
\label{b1:subsec:C10401}

本节固定 Banach 空间 \(X\) 及紧算子 \(K:X\rightarrow X\)，记 \(T=I-K\)。
由于紧算子在有界复合及线性组合下稳定，每个正整数 \(m\) 都有
\[
 T^m=I-K_m,\quad
 K_m=\sum_{j=1}^m(-1)^{j+1}\binom mjK^j,
\]
其中 \(K_m\) 仍紧。这个观察使下面关于 \(T\) 的结论也能用于它的所有正整数次幂。

\begin{lemma}\label{b1:lem:C1-compact-perturbation-range}
算子 \(T=I-K\) 的核有限维，值域闭；并且存在 \(C>0\)，使
\[
 \operatorname{dist}(x,\ker T)\leq C\|Tx\|\quad(x\in X).
\]
\end{lemma}
\begin{proof}
在闭子空间 \(N=\ker T\) 上，\(Kx=x\)。因此 \(N\) 的闭单位球在 \(X\) 中相对紧，
又因 \(N\) 闭而在 \(N\) 中紧。Riesz 引理的推论给出 \(N\) 有限维。

若所述估计不成立，对每个 \(n\) 可取 \(z_n\)，使
\(\operatorname{dist}(z_n,N)>n\|Tz_n\|\)。将 \(z_n\) 除以它到 \(N\) 的距离，
再减去一个 \(N\) 中的向量，可得到 \(x_n\)，满足
\[
 \operatorname{dist}(x_n,N)=1,\quad \|x_n\|\leq2,\quad \|Tx_n\|<1/n.
\]
减去核向量不改变 \(Tx_n\)，选取范数至多 \(2\) 的代表只用下确界定义。
紧性给出子列 \(Kx_{n_j}\rightarrow v\)，因而
\(x_{n_j}=Tx_{n_j}+Kx_{n_j}\rightarrow v\)。
连续性给出 \(Tv=0\)，但距离函数连续，所以
\(\operatorname{dist}(v,N)=1\)，矛盾。估计成立，再由%
\cref{b1:thm:C1-closed-range-estimate}得值域闭。
\end{proof}

若 \(T\) 单射，估计就控制整个向量。若 \(T\) 不单射，它仍控制与齐次解无关的部分。
证明没有依赖核的某个事先选好的补空间；先在商空间意义下归一化，
再借紧性得到极限，避免了“原像列有界”这一通常并不成立的假设。

\subsection{单射与满射为何等价}
\label{b1:subsec:C10402}

\begin{theorem}[Fredholm，单射与满射]\label{b1:thm:C1-fredholm-bijective}
对 \(T=I-K\)，下列条件等价：
\begin{equivalences}
\item\label{b1:item:C1-fredholm-injective}\(T\) 单射；
\item\label{b1:item:C1-fredholm-surjective}\(T\) 满射；
\item\label{b1:item:C1-fredholm-invertible}\(T\) 有有界双侧逆。
\end{equivalences}
\end{theorem}
\begin{proof}
先设 \(T\) 单射但不满射。记 \(R_n=T^nX\)，其中 \(R_0=X\)。
这些子空间闭，因为每个 \(T^n\) 都是恒等算子的紧扰动。
它们还严格递减：若 \(R_n=R_{n+1}\)，则对任意 \(x\in X\)，存在 \(y\) 使
\(T^nx=T^{n+1}y\)；由 \(T^n\) 单射，得到 \(x=Ty\)，与不满射矛盾。

对真闭子空间 \(R_n\subset R_{n-1}\) 应用 Riesz 引理，选取
\[
 x_n\in R_{n-1},\quad\|x_n\|=1,\quad\operatorname{dist}(x_n,R_n)>1/2.
\]
当 \(m>n\) 时，\(x_m\in R_n\)。又因为 \(K=I-T\) 与 \(T\) 交换，
\(K\) 保持每个 \(R_n\)，所以 \(Tx_n+Kx_m\in R_n\)。于是
\[
 \|Kx_n-Kx_m\|=\|x_n-(Tx_n+Kx_m)\|>1/2.
\]
紧性不允许有界列的像保持这种间距，矛盾。

再设 \(T\) 满射但不单射。记 \(N_n=\ker T^n\)，\(N_0=\{0\}\)。
从 \(0\ne z_1\in\ker T\) 出发，利用满射性逐次选取 \(Tz_{n+1}=z_n\)。
于是 \(z_n\in N_n\setminus N_{n-1}\)，各个闭子空间 \(N_n\) 严格递增。
仍由 Riesz 引理，可选 \(x_n\in N_n\)，使
\(\|x_n\|=1\) 且 \(\operatorname{dist}(x_n,N_{n-1})>1/2\)。
对 \(m<n\)，有 \(Tx_n\in N_{n-1}\)、\(Kx_m\in N_{n-1}\)，故
\[
 \|Kx_n-Kx_m\|=\|x_n-(Tx_n+Kx_m)\|>1/2,
\]
再次与紧性矛盾。单射与满射因此等价；双射的有界逆由开映射定理保证。
\end{proof}

在上述证明中，每次违背可逆性都会产生一条严格变化的子空间链。
Riesz 引理从链中选出彼此分开的方向，紧性再排除这样的无穷过程。
这种方法不依赖对角化，所以也适用于没有标准正交特征向量基的紧算子。

\subsection{有限个相容条件}
\label{b1:subsec:C10403}

\begin{definition}\label{b1:def:C1-fredholm-operator}
设 \(X,Y\) 为 Banach 空间，\(A\in\mathcal B(X,Y)\)。
若 \(\operatorname{ran}A\) 闭，且 \(\ker A\) 与商空间 \(Y/\operatorname{ran}A\)
均有限维，则称 \(A\) 为\term[fredholm-suanzi]{Fredholm 算子}[Fredholm operator]。
商空间 \(Y/\operatorname{ran}A\) 称为其\term[yuhe]{余核}[cokernel]，规定
\[
 \operatorname{ind}A=\dim\ker A-\dim(Y/\operatorname{ran}A),
\]
称为 \(A\) 的\term[fredholm-zhibiao]{Fredholm 指标}[Fredholm index]。
\end{definition}
\symidx{\(\operatorname{ind}A\)：Fredholm 算子的指标}

值域闭时，商空间对偶识别给出
\((Y/\operatorname{ran}A)^*=(\operatorname{ran}A)^\circ=\ker A'\)。
因此有限余维可由 \(\dim\ker A'<\infty\) 检测，且两维数相等。
这里还用到一个简单事实：若赋范空间的连续对偶有限维，则原空间也有限维且维数相同。
因为连续对偶分离点，自然映射把原空间单射到其有限维二次对偶；随后使用有限维对偶的维数公式。

\begin{theorem}[Fredholm，可解性与维数]\label{b1:thm:C1-fredholm-alternative}
设 \(K:X\rightarrow X\) 紧，\(T=I-K\)。则 \(T\) 为指标零的 Fredholm 算子：
\[
 \dim\ker T=\dim\ker T'<\infty,\quad
 \operatorname{ran}T={}^{\circ}(\ker T').
\]
因此恰有以下两种情形之一：
\begin{enumerate}
\item\label{b1:item:C1-fredholm-unique}核为零，且对每个 \(y\in X\)，方程 \(Tx=y\) 都有唯一解；
\item\label{b1:item:C1-fredholm-compatible}核非零且有限维；方程 \(Tx=y\) 有解，当且仅当
\(g(y)=0\) 对所有 \(g\in\ker T'\) 成立。有解时全部解为 \(x_0+\ker T\)。
\end{enumerate}
对偶方程 \(T'g=f\) 的可解性条件同样为 \(f(x)=0\) 对所有 \(x\in\ker T\) 成立。
\end{theorem}
\begin{proof}
由 Schauder 定理，\(K'\) 紧。所以%
\cref{b1:lem:C1-compact-perturbation-range}分别用于 \(T\) 与 \(T'=I-K'\)，
给出二者的有限维核及闭值域。Banach 闭值域定理给出两组可解性条件，
上面的商空间对偶识别则说明 \(T\) 具有有限维余核。
尚须证明两维数相同。

记 \(n=\dim\ker T\)，\(m=\dim(X/\operatorname{ran}T)\)，\(r=\min(n,m)\)。
在核中取基 \(z_1,\ldots,z_n\)，把其坐标泛函延拓为 \(f_1,\ldots,f_n\in X^*\)，
使 \(f_i(z_j)=\delta_{ij}\)。在 \(X\) 中选取 \(y_1,\ldots,y_r\)，使它们在
\(X/\operatorname{ran}T\) 中的像线性无关。定义有限秩算子
\[
 Fx=\sum_{i=1}^r f_i(x)y_i,\quad S=T+F.
\]
它仍是恒等算子的紧扰动，因为 \(S=I-(K-F)\)。

若 \(Sx=0\)，在商空间中取像，得到所有 \(f_i(x)=0\)，\(1\leq i\leq r\)；
于是 \(Tx=0\)。反过来，核中前 \(r\) 个坐标为零的向量都被 \(S\) 消去。因此
\(\dim\ker S=n-r\)。
另一方面，
\[
 \operatorname{ran}S=\operatorname{ran}T+\operatorname{span}\{y_1,\ldots,y_r\}.
\]
一个包含关系直接来自定义。反向关系中，\(Sz_i=y_i\)；
而对任意 \(x\)，向量 \(w=x-\sum_{i=1}^r f_i(x)z_i\) 满足 \(Sw=Tx\)。
所以等式成立，余维为 \(m-r\)。

按 \(r\) 的选择，\(S\) 的核与余核至少有一个为零。
由\cref{b1:thm:C1-fredholm-bijective}，单射与满射等价，所以另一个也为零。
故 \(n=r=m\)。若这个公共维数为零，可逆性给出唯一可解；
若非零，可解性由闭值域定理刻画，任意两个解的差属于核。
\end{proof}

这个可解性结论在此记为\term[fredholm-dingli]{Fredholm 定理}[Fredholm alternative]。
在实际计算中，先求一组 \(\ker T'\) 的基 \(g_1,\ldots,g_n\)，
条件便缩减为 \(n\) 个标量等式 \(g_j(y)=0\)。相容条件的数目与解的不唯一方向数相等，
而不是与所选近似空间的维数相等。

指标零是紧扰动恒等算子的特殊性质，不是 Fredholm 算子的定义。
例如右移算子在 \(\ell^2\) 上单射且值域余维为 \(1\)，所以指标为 \(-1\)；
其伴随左移满射而核一维，指标为 \(1\)。这两个算子都不能表示为恒等算子加一个紧算子。

\subsection{非零谱点的有限维性}
\label{b1:subsec:C10404}

以下仅在复 Banach 空间上谈谱，并沿用\chapref{b1:ch:08}的记号 \(\sigma(K)\)：
\(\lambda\) 属于谱是指 \(\lambda I-K\) 没有有界双侧逆。

\begin{corollary}[紧算子的非零谱]\label{b1:cor:C1-compact-nonzero-spectrum}
若 \(K\) 紧，则每个非零谱点都是特征值，对应特征空间有限维。
对每个 \(\varepsilon>0\)，绝对值至少为 \(\varepsilon\) 的不同特征值只有有限个。
因此非零谱点至多可数，且除零以外没有聚点。
\end{corollary}
\begin{proof}
对 \(\lambda\ne0\)，将 Fredholm 定理用于 \(I-\lambda^{-1}K\)。
若没有非零齐次解，它便可逆，所以非零谱点必是特征值；
核有限维给出特征空间的结论。

若有无穷多个不同特征值 \(\lambda_n\)，且 \(|\lambda_n|\geq\varepsilon\)，
各选特征向量 \(v_n\ne0\)，令 \(M_n=\operatorname{span}\{v_1,\ldots,v_n\}\)。
不同特征值的特征向量线性无关：若有最短非平凡线性关系，对关系作用
\(K-\lambda_nI\) 就得到更短的关系。故 \(M_{n-1}\) 是 \(M_n\) 的真子空间。
取 \(x_n\in M_n\)，使 \(\|x_n\|=1\) 且到 \(M_{n-1}\) 的距离大于 \(1/2\)。
由于 \((K-\lambda_nI)M_n\subset M_{n-1}\)，当 \(m<n\) 时，
\[
 \frac{Kx_n}{\lambda_n}-\frac{Kx_m}{\lambda_m}
 =x_n+w,\quad w\in M_{n-1},
\]
两项之差的范数大于 \(1/2\)。但 \(|\lambda_n|\leq\|K\|\)，可先抽取
\(\lambda_n\) 的收敛子列，极限的模至少为 \(\varepsilon\)，再利用 \(K\) 紧抽取
\(Kx_n\) 的收敛子列。所得 \(Kx_n/\lambda_n\) 收敛，矛盾。
最后对 \(\varepsilon=1,1/2,\ldots\) 取并，得到可数性与聚点结论。
\end{proof}

非零特征空间有限维，并不意味着这些空间张成全空间。没有自伴性时，
有限维部分可以含 Jordan 型行为；甚至可能没有任何非零特征值。
本附录习题中的加权移位算子就是后一种情形。下一节恢复自伴性，才能进一步把可解性写成正交坐标条件。

% !TeX root = ../../Book1.tex
\section{紧自伴谱分解与方程求解}
\label{b1:sec:C105}

上一节没有使用内积，因此得到了广泛的可解性结论，却没有给出解的坐标。
对于紧自伴算子，\cref{b1:thm:compact-self-adjoint-spectral}已经建立完整的正交分解。
这里把它用于三件事：写出相容条件，估计截断误差，以及区分两类性质很不相同的方程。
谱定理的证明不再重复，重点是从分解到求解之间的步骤。

\subsection{零特征空间与非零坐标}
\label{b1:subsec:C10501}

设 \(H\) 为实或复 Hilbert 空间，\(K:H\rightarrow H\) 紧且自伴。
按谱定理选取有限个或可数个标准正交向量 \(e_j\)，对应非零实特征值 \(\lambda_j\)，并记
\[
 H_0=\ker K,\quad
 x=x_0+\sum_jx_je_j,\quad
 x_j=\langle x,e_j\rangle,\quad x_0=P_{H_0}x.
\]
全部级数均在 \(H\) 中按范数收敛。若非零特征值只有有限个，以下的和与条件
只遍历实际存在的指标；若 \(K=0\)，这些和为空和。
核 \(H_0\) 不要求有限维或可分，不能把它遗漏后声称 \(\{e_j\}\) 是 \(H\) 的基。

\begin{proposition}\label{b1:prop:C1-spectral-solvability}
给定标量 \(\alpha\)，令 \(E_\alpha=\ker(I-\alpha K)\)。
方程
\[
 (I-\alpha K)x=y
\]
有解，当且仅当 \(y\perp E_\alpha\)。满足 \(x\perp E_\alpha\) 的解唯一，且为
\[
 x=y_0+
 \sum_{1-\alpha\lambda_j\ne0}
       \frac{y_j}{1-\alpha\lambda_j}e_j,
 \quad y_j=\langle y,e_j\rangle,\quad y_0=P_{H_0}y.
\]
所有解由这个解加上 \(E_\alpha\) 中的任意向量得到。
\end{proposition}
\begin{proof}
在 \(H_0\) 上，\(I-\alpha K\) 为恒等算子；在 \(e_j\) 方向上，它为乘
\(1-\alpha\lambda_j\)。所以方程要求 \(x_0=y_0\)，以及
\[
 (1-\alpha\lambda_j)x_j=y_j.
\]
当乘子为零时必须有 \(y_j=0\)，这些方向恰好张成 \(E_\alpha\)。
若 \(\alpha=0\)，这个空间为零；若 \(\alpha\ne0\)，它至多是一个非零特征值
对应的有限维特征空间。

还须验证候选解的级数收敛。非零分母存在统一正下界 \(\delta_\alpha\)：
无穷列情形中 \(\lambda_j\to0\)，故充分大的 \(j\) 满足
\(|1-\alpha\lambda_j|\geq1/2\)；剩下有限多个非零分母的模也有正最小值。
有限列情形直接取有限个正数的最小值，没有此类指标时取 \(\delta_\alpha=1\)。
由 Bessel 不等式，
\[
 \sum_{1-\alpha\lambda_j\ne0}
 \frac{|y_j|^2}{|1-\alpha\lambda_j|^2}
 \leq\delta_\alpha^{-2}\|y\|^2<\infty.
\]
级数因而收敛。对部分和作用连续算子并取极限，得到原方程。
其在 \(E_\alpha\) 上的分量为零；任意两个解之差属于 \(E_\alpha\)，所以唯一性成立。
\end{proof}

在复空间中，\(\alpha\) 可以非实。此时 \(I-\alpha K\) 不一定自伴，但其核与伴随的核相同：
伴随为 \(I-\overline\alpha K\)，实特征值使两个算子可能出现的零乘子完全一致。
因此这里的正交条件与前一节的 Fredholm 条件相符，不是忽略了复共轭。

\ProseParagraphBreak

当参数接近某个 \(\lambda_j^{-1}\) 时，相应方向的分母变小，解会对右端误差更敏感。
在参数恰等于这个值时，该方向不再由方程确定，同时右端在该方向上的分量必须为零。
这分别是接近共振和精确共振的两种情况。

\subsection{有限秩截断的稳定性}
\label{b1:subsec:C10502}

设 \(I-\alpha K\) 可逆，记
\[
 K_Nx=\sum_{j=1}^N\lambda_j\langle x,e_j\rangle e_j,\quad
 \varepsilon_N=\|K-K_N\|.
\]
有限秩情形取足够大的 \(N\) 后令 \(K_N=K\)。截断的对象是算子，不是先把整个右端扔进
有限维空间。因此 \(I-\alpha K_N\) 在所选特征向量的正交补上仍为恒等算子。

\begin{proposition}\label{b1:prop:C1-spectral-truncation}
令 \(A=I-\alpha K\)，\(A_N=I-\alpha K_N\)，并设
\[
 q_N=|\alpha|\cdot\|A^{-1}\|\varepsilon_N<1.
\]
则 \(A_N\) 可逆。若 \(Ax=y\)、\(A_Nx_N=y\)，则
\[
 \|x_N-x\|\leq
 \frac{|\alpha|\cdot\|A^{-1}\|^2\varepsilon_N}{1-q_N}\|y\|.
\]
特别地，\(x_N\to x\)，且此收敛对有界的右端一致。
\end{proposition}
\begin{proof}
有
\[
 A_N=A\bigl(I+A^{-1}\alpha(K-K_N)\bigr).
\]
括号中的扰动范数不超过 \(q_N\)。若 \(\|B\|<1\)，则
\(\sum_{m\geq0}(-B)^m\) 在算子范数下收敛，与 \(I+B\) 相乘后由有限几何级数公式
及余项趋零得到双侧逆；其范数不超过 \((1-\|B\|)^{-1}\)。
因此
\[
 \|A_N^{-1}\|\leq\frac{\|A^{-1}\|}{1-q_N}.
\]
再由 \(A_N(x_N-x)=\alpha(K_N-K)x\)，得到所述误差估计。
谱定理给出 \(\varepsilon_N\to0\)，所以充分大的 \(N\) 满足假设。
\end{proof}

这一估计同时说明了截断误差和逆算子范数的作用。仅知道 \(K_N\to K\)，
却不检查极限方程是否可逆，不能保证离散方程的逆一致有界。
若固定在共振参数，必须先施加相容条件、去掉核，再在其正交补上进行相同论证。

\subsection{直接求逆与小特征值}
\label{b1:subsec:C10503}

方程 \(Kx=y\) 没有恒等项。它的分母是 \(\lambda_j\)，无穷多个小特征值
会改变可解性条件，不能照搬上一小节中的统一下界。

\begin{proposition}\label{b1:prop:C1-spectral-range}
方程 \(Kx=y\) 有解，当且仅当
\[
 y\perp\ker K,\quad \sum_j\frac{|y_j|^2}{|\lambda_j|^2}<\infty.
\]
有解时，唯一的最小范数解为
\[
 x^\dagger=\sum_j\frac{y_j}{\lambda_j}e_j.
\]
此外，紧自伴算子的值域闭，当且仅当其秩有限。
\end{proposition}
\begin{proof}
若 \(Kx=y\)，则 \(y\perp\ker K\)，且 \(x_j=y_j/\lambda_j\)；
Bessel 不等式给出所述可求和条件。反过来，该条件保证候选级数收敛，
对部分和作用 \(K\) 后取极限就得到 \(Kx^\dagger=y\)。
它属于 \((\ker K)^\perp\)，由正交分解得到最小范数性。

有限秩时值域有限维，因而闭。若秩无限，则非零特征向量有无穷多个，
\(\lambda_j\to0\)。每个 \(e_j\) 到 \(\ker K\) 的距离等于 \(1\)，
但 \(\|Ke_j\|=|\lambda_j|\to0\)，违反闭值域的估计判别。
\end{proof}

例如 \(Ke_j=j^{-1}e_j\) 时，右端 \(y=\sum_{j\geq1}j^{-1}e_j\) 属于
\(\ell^2\)，也与核正交，却不在值域中，因为候选原像的每个坐标都为 \(1\)。
取前 \(N\) 个坐标所得右端趋于 \(y\)，其解的范数却等于 \(\sqrt N\)。
正交相容只排除了核的障碍，小特征值另带来无穷维的可求和要求。

\ProseParagraphBreak

对确实可解的右端，截断
\(\sum_{j=1}^N(y_j/\lambda_j)e_j\) 总趋于最小范数解；
但若观测右端有微小误差，除以小 \(\lambda_j\) 会放大它。
习题进一步给出一个带正参数的稳定近似，使读者能直接比较“极限存在”
与“每个近似对数据连续”这两个要求。

求解线性方程时，值域闭性与有限维缺陷各自承担不同作用，见\cref{b1:tab:visual-t17}。
% Source fingerprint: f845e21e27a6874ee002bb345436ca13545a1e22a1fa5ebbe9c9e09117f3263b
% T17: raw TeX cells preserved; no numeric highlighting.
\begin{table}[htbp]
\centering\small\renewcommand{\arraystretch}{1.18}\setlength{\tabcolsep}{4pt}
\caption[Fredholm 可解性的检查项目]{Fredholm 可解性的检查项目。核有限维、值域闭和指标零是三个不同结论。一般 Banach 空间先用对偶配对，不能无条件借用正交语言。}
\label{b1:tab:visual-t17}
\begin{tabularx}{\linewidth}{@{}>{\raggedright\arraybackslash}p{3.1cm}>{\raggedright\arraybackslash}p{4.7cm}>{\raggedright\arraybackslash}X@{}}
\toprule
算子情形 & 已获得的结构 & 方程的相容条件 \\
\midrule
一般有界线性算子 & 核闭；值域不一定闭 & 不能因存在近似解就断言右端属于值域 \\
闭值域算子 & \(X/\ker A\) 上有由 \(\|Ax\|\) 控制的估计 & 闭值域定理把值域识别为相应对偶核的零化子 \\
\(I-K\)，\(K\) 紧 & 有限维核、闭值域、有限维余核，指标为零 & 单射等价于满射；不要求 \(K\) 自伴或可以对角化 \\
Hilbert 空间中的闭值域 & 可用伴随和正交补表示 & \(Ax=y\) 可解当且仅当 \(y\perp\ker A^*\) \\
紧自伴谱方法 & 沿正交特征向量逐分量处理 & 零特征值处有相容条件；小非零特征值还要求解系数可平方求和 \\
\bottomrule
\end{tabularx}
\end{table}


% !TeX root = ../../Book1.tex
\section{凸分离的进一步形式}
\label{b1:sec:C106}

\chapref{b1:ch:08}的几何 Hahn--Banach 定理把凸集与其外的一点分开，
本附录的闭值域证明已经用它把对偶估计转成有界的近似原像。
本节沿另外两条方向继续：把一点换成紧凸集，以及把“分开”定量化为距离。
这些形式只需赋范空间，不预先要求完备性或内积。

\subsection{紧凸集与闭凸集的统一间隙}
\label{b1:subsec:C10601}

\begin{theorem}[Hahn--Banach，紧集与闭集的凸分离]
\label{b1:thm:C1-compact-closed-separation}
设 \(A,C\) 是实或复赋范空间 \(X\) 中不交的非空凸集，\(A\) 紧、\(C\) 闭。
则存在 \(f\in X^*\) 及实数 \(a<b\)，使
\[
 \sup_{c\in C}\operatorname{Re}f(c)\leq a
 <b\leq\inf_{u\in A}\operatorname{Re}f(u).
\]
\end{theorem}
\begin{proof}
差集 \(D=C-A\) 凸，且不含零。先证明 \(D\) 闭：
若 \(c_n-u_n\to z\)，由 \(A\) 紧可取子列使 \(u_{n_j}\to u\in A\)，
于是 \(c_{n_j}\to z+u\in C\)，故 \(z\in D\)。
赋范空间为度量空间，这个序列判据足以给出闭性。

对 \(D\) 和零点应用\cref{b1:thm:geometric-hahn-banach}，可取非零 \(f\in X^*\) 与
\(\gamma<0\)，使
\[
 \operatorname{Re}f(c-u)\leq\gamma\quad(c\in C,\ u\in A).
\]
连续函数 \(\operatorname{Re}f\) 在紧集 \(A\) 上有有限最小值 \(b\)。
上式给出 \(\sup_C\operatorname{Re}f\leq b+\gamma\)。
令 \(a=b+\gamma\)，即得严格的统一间隙。
\end{proof}

紧性用在差集的闭性中，而不是只用来保证 \(A\) 有界。
若只要求两个集合闭，差集未必闭，统一间隙也可能消失。
这与\chapref{b1:ch:08}给出的指数曲线上方区域和横轴的例子一致。
另一方面，任何这样的泛函都给出
\[
 \|u-c\|\geq\frac{b-a}{\|f\|}\quad(u\in A,\ c\in C),
\]
所以统一分离必定伴随正距离。

\subsection{用连续泛函计算距离}
\label{b1:subsec:C10602}

对非空凸集 \(C\subset X\)，给定 \(f\in X^*\) 后，上确界
\(\sup_{c\in C}\operatorname{Re}f(c)\) 可以是 \(+\infty\)。
下面的公式允许这一情形，相应的差值记为 \(-\infty\)；零泛函始终给出值零。

\begin{theorem}\label{b1:thm:C1-distance-dual}
若 \(C\) 是非空闭凸集，则对每个 \(x\in X\)，
\[
 \operatorname{dist}(x,C)=
 \sup_{\|f\|\leq1}
 \left(\operatorname{Re}f(x)-\sup_{c\in C}\operatorname{Re}f(c)\right).
\]
若 \(x\notin C\)，右边的上确界可由某个范数为 \(1\) 的泛函取得。
这个结论不要求 \(C\) 中存在离 \(x\) 最近的点。
\end{theorem}
\begin{proof}
对任意 \(\|f\|\leq1\) 及 \(c\in C\)，有
\[
 \operatorname{Re}f(x)-\sup_C\operatorname{Re}f
 \leq \operatorname{Re}f(x-c)\leq\|x-c\|.
\]
先取距离下确界，再对泛函取上确界，得右边不大于左边。
若 \(x\in C\)，两边都等于零。

设 \(d=\operatorname{dist}(x,C)>0\)。开球 \(B(x,d)\) 与 \(C\) 是不交非空凸集，
由开集一侧的分离定理，改变泛函的符号后可取非零 \(f\in X^*\)，使
\[
 \sup_C\operatorname{Re}f\leq
 \inf_{z\in B(x,d)}\operatorname{Re}f(z)
 =\operatorname{Re}f(x)-d\|f\|.
\]
最后一个等式来自
\(\inf_{\|v\|<d}\operatorname{Re}f(v)=-d\|f\|\)：
实情形可变号，复情形可旋转标量相位，使模接近范数的值变为负实数。
除以 \(\|f\|\) 即得一个单位泛函，使右边至少为 \(d\)。
结合已经证明的反向不等式，等号及取得性成立。
\end{proof}

该公式把寻找近点的几何问题变成寻找测试泛函的问题。
对于闭子空间，有限的上确界迫使泛函消去整个子空间，便回到%
\cref{b1:thm:quotient-banach}与\cref{b1:prop:quotient-dual}给出的商范数与商空间对偶公式。
一般凸集不具备正负缩放的封闭性，因此这里必须保留实部和整个上确界。

\begin{corollary}\label{b1:cor:C1-best-approximation}
设 \(C\) 为非空闭凸集，\(x\notin C\)，\(c_0\in C\)。
则 \(c_0\) 是离 \(x\) 最近的点，当且仅当存在 \(\|f\|=1\) 的连续线性泛函，使
\[
 \operatorname{Re}f(x-c_0)=\|x-c_0\|,\quad
 \operatorname{Re}f(c-c_0)\leq0\quad(c\in C).
\]
\end{corollary}
\begin{proof}
若 \(c_0\) 最近，在距离公式中取到上确界的单位泛函满足
\[
 d=\operatorname{Re}f(x)-\sup_C\operatorname{Re}f
 \leq\operatorname{Re}f(x-c_0)\leq\|x-c_0\|=d.
\]
所有不等号均为等号，得到两个条件。
反过来，若条件成立，则对每个 \(c\in C\)，
\[
 \|x-c\|\geq\operatorname{Re}f(x-c)
 =\|x-c_0\|-\operatorname{Re}f(c-c_0)\geq\|x-c_0\|.
\]
所以 \(c_0\) 最近。
\end{proof}

在 Hilbert 空间中可取
\(f(v)=\langle v,x-c_0\rangle/\|x-c_0\|\)，于是第二个条件成为
\[
 \operatorname{Re}\langle c-c_0,x-c_0\rangle\leq0\quad(c\in C).
\]
这就是凸集最近点的变分条件。若 \(C\) 为线性子空间，沿正负和复数倍方向变化后，
它化为通常的正交条件。

\subsection{边界处的线性支撑及其限制}
\label{b1:subsec:C10603}

距离公式处理的是集外点。若点已在凸集边界上，则正距离消失；
仍想找到非零泛函在该点取得最大值，需要另看集合是否具有内部。

\begin{proposition}\label{b1:prop:C1-boundary-support}
设 \(C\subset X\) 为闭凸集，内部非空。对每个边界点 \(c_0\in\partial C\)，
存在非零 \(f\in X^*\)，使
\[
 \operatorname{Re}f(c)\leq\operatorname{Re}f(c_0)\quad(c\in C).
\]
\end{proposition}
\begin{proof}
内部 \(U=\operatorname{int}C\) 为开凸集，且不含 \(c_0\)。
开凸集分离定理给出非零 \(f\)，使
\(\operatorname{Re}f(u)\leq\operatorname{Re}f(c_0)\) 对每个 \(u\in U\) 成立。
只须证明 \(\overline U=C\)。
取 \(u_0\in U\) 和 \(r>0\)，使 \(B(u_0,r)\subset C\)。
对任意 \(c\in C\) 与 \(0<t<1\)，凸性给出
\[
 B\bigl((1-t)c+tu_0,tr\bigr)
 =(1-t)c+tB(u_0,r)\subset C.
\]
故 \((1-t)c+tu_0\in U\)，当 \(t\downarrow0\) 时趋于 \(c\)。
由 \(C\) 闭及 \(f\) 连续，结论成立。
\end{proof}

内部条件在无穷维中不能随意删除。以下反例只针对实空间；
它已经足以说明一般无穷维赋范空间中不能删除内部非空假设。
具体地，在实 \(\ell^2\) 中取闭凸集
\[
 C=\{\,x:x_j\geq0\ \text{对每个 }j\,\}.
\]
它的内部为空：任一 \(x\in C\) 的坐标趋零，因此任意小的球中都能把某个远处坐标
改成负值。取 \(c_0=(2^{-j})_{j\geq1}\)，这是一个边界点。
若 \(f\) 满足上面的支撑不等式，则对每个固定 \(j\)，向量
\(c_0\pm te_j\) 在充分小的正 \(t\) 下都属于 \(C\)，从而 \(f(e_j)=0\)。
有限支撑向量稠密，连续性迫使 \(f=0\)。这个点不存在所要求的非零泛函。

分离集外点、统一分离两个集合、在边界点作线性支撑，是三个不同的问题。
它们都来自 Hahn--Banach 定理，却分别依赖空球、紧性和内部条件。
使用时先确认需要哪一种结论，往往比记住一个笼统的“分离定理”名称更有用。


\begin{chaptersummary}
方程 \(Tx=y\) 的不唯一性由核描述，近似可解性由对偶核描述；
闭值域估计使近似可解性与精确可解性一致，并在商空间中控制解对右端的依赖。
这些论断对一般 Banach 空间成立。Hilbert 空间的正交补额外提供唯一的最小范数代表，
不能在没有内积时直接照搬。

紧性把可能失控的方向压缩为有限维障碍。对于 \(I-K\)，
核有限维、值域闭和指标为零经过不同的证明步骤得到；
右移与左移的例子说明，最后一个性质不能从前两个直接推出。
自伴性再把可解性写成正交坐标条件。
恒等项防止高阶坐标的分母趋零，而直接求解 \(Kx=y\) 时，小特征值会造成额外的可求和条件。

凸分离则把几何约束写成线性测试。集外一点总有距离公式，
紧凸集与不交闭凸集之间有统一间隙，内部非空的闭凸集在边界处有非零线性支撑。
三种结论各自使用不同的几何假设；在具体的方程或极小化问题中，
应把这些假设与所需结论一并带入。
\end{chaptersummary}

\BookExercises

% 八题均为自拟；逐题注释列出训练增量。
% 不复述\chapref{b1:ch:08}商空间、c0对偶及\chapref{b1:ch:09}双零化子、范数不取得的既有习题。

\begin{exercise}\label{b1:ex:C1-nonclosed-sum}
在实或复 \(H=\ell^2\) 上令 \(De_j=j^{-1}e_j\)，并在 \(H\oplus H\) 中取
\[
 M=\{\,(x,Dx):x\in H\,\},\quad N=H\oplus\{0\}.
\]
证明 \(M,N\) 都是闭线性子空间，\(M\cap N=\{0\}\)，但 \(M+N\) 不闭。
求 \(\overline{M+N}\)，并构造一列和空间中的向量，其和收敛而 \(M\) 分量的范数发散。
\end{exercise}
% 来源：自拟；把非闭值域变成两个闭子空间之和不闭，增加分解稳定性的解释。
\begin{solution}
若 \((x_n,Dx_n)\to(x,y)\)，由 \(D\) 连续得到 \(y=Dx\)，所以图像 \(M\) 闭。
空间 \(N\) 是第二坐标投影的核，也闭。由于 \(D\) 单射，交集为零。
直接计算给出
\[
 M+N=H\oplus\operatorname{ran}D.
\]
值域包含所有有限支撑序列，故稠密；但
\(y=\sum_{j\geq1}j^{-1}e_j\) 不在值域中，其唯一可能原像是常数序列 \(1\)。
因此和空间不闭，而其闭包等于 \(H\oplus H\)。

令 \(u_n=\sum_{j=1}^ne_j\)，则
\[
 (u_n,Du_n)+(-u_n,0)=(0,Du_n)\longrightarrow(0,y).
\]
这是由于交集为零而唯一的 \(M+N\) 分解；
其中 \(M\) 分量的范数至少为 \(\|u_n\|=\sqrt n\)。
两个闭子空间分别具有完备性，却没有给它们的和提供有界的分量提取。
\end{solution}

\begin{exercise}\label{b1:ex:C1-rank-one-fredholm}
设 \(X\) 为 Banach 空间，\(u\in X\setminus\{0\}\)、\(f\in X^*\setminus\{0\}\)，
令 \(Kx=f(x)u\)。对任意标量 \(\alpha\)，分别在
\(1-\alpha f(u)\ne0\) 与 \(1-\alpha f(u)=0\) 时求解
\[
 (I-\alpha K)x=y.
\]
写出 \(K'\)，并在第二种情形直接核对对偶相容条件。
\end{exercise}
% 来源：自拟；在复Banach配对中显式计算秩一算子，避免把对偶误写成Hilbert伴随。
\begin{solution}
恒有 \(K^2=f(u)K\)。当 \(1-\alpha f(u)\ne0\) 时，直接相乘得到
\[
 (I-\alpha K)^{-1}
 =I+\frac{\alpha}{1-\alpha f(u)}K,
\]
所以唯一解为
\[
 x=y+\frac{\alpha f(y)}{1-\alpha f(u)}u.
\]
这包括 \(\alpha=0\) 或 \(f(u)=0\) 的情况。

当 \(\alpha f(u)=1\) 时，对方程作用 \(f\) 得 \(f(y)=0\)。
若该条件满足，则 \(x=y\) 是一解；齐次方程要求 \(x=\alpha f(x)u\)，
故核恰为 \(\mathbb Ku\)，所有解为 \(y+tu\)。
对任意 \(g\in X^*\)，
\[
 K'g=g(u)f.
\]
因此 \((I-\alpha K')g=0\) 恰好等价于 \(g\in\mathbb Kf\)。
对偶核消去右端的条件就是 \(f(y)=0\)，与直接计算一致。
这里没有对 \(\alpha\) 或 \(f(u)\) 加复共轭，因为连续对偶由线性泛函组成。
\end{solution}

\begin{exercise}\label{b1:ex:C1-shift-index}
在 \(\ell^2(\mathbb N)\) 上令
\[
 R(x_1,x_2,\ldots)=(0,x_1,x_2,\ldots),\quad
 L(x_1,x_2,\ldots)=(x_2,x_3,\ldots).
\]
给定非负整数 \(m,n\)，求 \(A=R^mL^n\) 的核、值域、伴随和指标。
判定 \(A\) 何时可以写成 \(I+F\)，其中 \(F\) 有限秩；
再判定 \(m\ne n\) 时它能否写成 \(I+K\)，其中 \(K\) 紧。
\end{exercise}
% 来源：自拟；把单次左右移推广到不同次数的复合，区分零指标与具体有限秩扰动。
\begin{solution}
算子 \(L^n\) 丢掉前 \(n\) 个坐标，\(R^m\) 再在前面补 \(m\) 个零，所以
\[
 \ker A=\operatorname{span}\{e_1,\ldots,e_n\},\quad
 \operatorname{ran}A=\{\,y:y_1=\cdots=y_m=0\,\}.
\]
当次数为零时相应限制为空。值域闭且余维为 \(m\)，故
\(\operatorname{ind}A=n-m\)。
由 \(R^*=L\)、\(L^*=R\)，有 \(A^*=R^nL^m\)，其核维数为 \(m\)。

若 \(m=n\)，则 \(A=I-P_n\)，其中 \(P_n\) 是前 \(n\) 个坐标的正交投影，
所以确为有限秩扰动。若 \(m\ne n\)，指标非零，
而恒等算子的紧扰动指标必须为零，因此不能写成 \(I+K\)，更不能写成有限秩扰动。
也可以直接看见非紧性：取充分稀疏的一列 \(j_k>n\)，使
\(\{j_k,j_k-n+m\}\) 互不相交，则 \((A-I)e_{j_k}\) 两两正交且范数均为 \(\sqrt2\)，
不存在收敛子列。
\end{solution}

\begin{exercise}\label{b1:ex:C1-weighted-shift}
在复 \(\ell^2(\mathbb N)\) 上定义 \(Ke_j=j^{-1}e_{j+1}\)。
证明 \(K\) 紧且没有非零特征值，进一步证明
\[
 \|K^r\|=\frac1{r!}\quad(r\geq1),\quad
 (I-\alpha K)^{-1}=\sum_{r=0}^{\infty}\alpha^rK^r
 \quad(\alpha\in\mathbb C).
\]
由此求 \(\sigma(K)\)，并说明非零特征向量不一定足以描述紧算子。
\end{exercise}
% 来源：自拟；提供非自伴且无非零特征值的紧算子，幂范数给出任意参数的收敛逆。
\begin{solution}
把 \(K\) 限制到前 \(N\) 个输入坐标得到有限秩算子 \(K_N\)，且
\(\|K-K_N\|=1/(N+1)\)。所以 \(K\) 紧。
若 \(Kx=\lambda x\) 且 \(\lambda\ne0\)，第一坐标给出 \(x_1=0\)，
以后逐个坐标给出 \(x_j=0\)。故没有非零特征值。

重复作用得到
\[
 K^re_j=\frac{(j-1)!}{(j+r-1)!}e_{j+r}.
\]
这些输出方向正交，系数随 \(j\) 递减，最大值在 \(j=1\) 处为 \(1/r!\)。
对任意 \(\alpha\)，算子级数由
\(\sum_r|\alpha|^r/r!<\infty\) 绝对控制，故在算子范数下收敛。
其前 \(N+1\) 项与 \(I-\alpha K\) 左右相乘，余项均为
\(-\alpha^{N+1}K^{N+1}\)，范数趋零，因此给出双侧逆。
所以每个 \(\lambda\ne0\) 都不在谱中；而 \(K\) 的像第一坐标恒为零，不满射，
故 \(0\in\sigma(K)\)。于是 \(\sigma(K)=\{0\}\)，但 \(K\ne0\)。
紧性没有提供自伴算子所具有的特征向量展开。
\end{solution}

\begin{exercise}\label{b1:ex:C1-spectral-regularization}
沿用\cref{b1:prop:C1-spectral-range}中的紧自伴算子 \(K\)。
对 \(\eta>0\)、\(y\in H\)，定义
\[
 x_\eta(y)=\sum_j\frac{\lambda_j}{\lambda_j^2+\eta}y_je_j.
\]
证明级数收敛，且这是 \((K^2+\eta I)x=Ky\) 的唯一解，并满足
\[
 \|x_\eta(y)-x_\eta(z)\|\leq\frac1{2\sqrt\eta}\|y-z\|.
\]
若 \(y=Kx^\dagger\)、\(x^\dagger\perp\ker K\)，证明 \(x_\eta(y)\to x^\dagger\)。
若还可写 \(x^\dagger=Kw\)，证明
\[
 \|x_\eta(y)-x^\dagger\|\leq\frac{\sqrt\eta}{2}\|w\|.
\]
最后给出右端误差不超过 \(\delta\) 时的总误差界。
\end{exercise}
% 来源：自拟；从谱坐标构造稳定近似并量化噪声与逼近误差，不引用尚未建立的正则化理论。
\begin{solution}
由 \(2\sqrt\eta\,|\lambda|\leq\lambda^2+\eta\)，每个乘子绝对值不超过
\(1/(2\sqrt\eta)\)。正交级数判据给出收敛及题述 Lipschitz 估计。
算子 \(K^2+\eta I\) 在 \(H_0\) 上乘 \(\eta\)，在 \(e_j\) 上乘
\(\lambda_j^2+\eta\)，均以 \(\eta\) 为正下界。
逐坐标可知候选向量满足方程，齐次方程各分量均为零，故解唯一。

若 \(y_j=\lambda_jx^\dagger_j\)，则
\[
 x_\eta(y)-x^\dagger
 =-\sum_j\frac{\eta}{\lambda_j^2+\eta}x^\dagger_je_j.
\]
每个系数乘子趋零且绝对值至多为 \(1\)。先使
\(\sum_{j>N}|x^\dagger_j|^2\) 任意小，再对有限个首项取 \(\eta\downarrow0\)，
得到范数收敛。
若 \(x^\dagger=Kw\)，误差的乘子改为
\(\eta\lambda_j/(\lambda_j^2+\eta)\)，其绝对值不超过 \(\sqrt\eta/2\)，
于是得到更精确的界。

若 \(\|y^\delta-y\|\leq\delta\)，在最后这个附加条件下，
\[
 \|x_\eta(y^\delta)-x^\dagger\|
 \leq\frac{\delta}{2\sqrt\eta}+\frac{\sqrt\eta}{2}\|w\|.
\]
例如令 \(\eta=\delta\) 并使 \(\delta\downarrow0\)，两项都趋零。
若固定 \(\delta>0\) 却一味使 \(\eta\) 趋零，第一项可能增大，不能只依据精确数据的极限选择参数。
\end{solution}

\begin{exercise}\label{b1:ex:C1-galerkin}
设可分 Hilbert 空间 \(H\) 有标准正交基 \((e_j)\)，\(P_N\) 为前 \(N\) 个坐标的正交投影。
设 \(K:H\rightarrow H\) 紧，不要求自伴，且 \(A=I-K\) 可逆。
证明
\[
 \|(I-P_N)K\|\longrightarrow0.
\]
由此证明：对充分大的 \(N\)，有限维方程
\[
 x_N-P_NKx_N=P_Ny,\quad x_N\in P_NH
\]
有唯一解；若 \(Ax=y\)，则存在与 \(N,y\) 无关的常数 \(C\)，使
\[
 \|x_N-x\|\leq C\|(I-P_N)x\|.
\]
\end{exercise}
% 来源：自拟；将自伴谱截断换成任意正交基投影，需用紧像有限网把逐点收敛提升为算子范数收敛。
\begin{solution}
投影满足 \(\|I-P_N\|\leq1\)，并对每个 \(z\in H\) 有 \(P_Nz\to z\)。
给定 \(\varepsilon>0\)，在紧集 \(\overline{K(B_H)}\) 中取半径
\(\varepsilon/2\) 的有限网 \(z_1,\ldots,z_m\)。
对充分大的 \(N\)，每个网点都满足
\(\|(I-P_N)z_i\|<\varepsilon/2\)。
先移到网点，再用投影范数控制误差，即得
\(\sup_{\|v\|\leq1}\|(I-P_N)Kv\|\leq\varepsilon\)。

令 \(A_N=I-P_NK\)。则
\[
 A_N=A\bigl(I+A^{-1}(I-P_N)K\bigr).
\]
扰动范数趋零，由几何级数构造逆，充分大时有
\(\|A_N^{-1}\|\leq2\|A^{-1}\|=:C\)。
全空间中的唯一解 \(x_N=A_N^{-1}P_Ny\) 自动满足
\((I-P_N)x_N=0\)，故确在 \(P_NH\) 内，给出题中的有限维解。
又因 \(Ax=y\)，
\[
 A_N(x_N-x)=P_Ny-x+P_NKx=-(I-P_N)x.
\]
应用 \(A_N^{-1}\) 得误差界。右边是解的正交截断误差，而不是预先假定的数值求解误差。
\end{solution}

\begin{exercise}\label{b1:ex:C1-dual-maximum-no-nearest}
\begin{enumerate}
\item\label{b1:item:C1-nonunique-nearest}
在 \(\mathbb R^2\) 的最大值范数下，求点 \((0,1)\) 到横轴的全部最近点。
\item\label{b1:item:C1-no-nearest}
在实 \(c_0\) 上令 \(f(x)=\sum_{j\geq1}2^{-j}x_j\)，取 \(C=\ker f\)、\(x=2e_1\)。
证明 \(\operatorname{dist}(x,C)=1\)，但 \(C\) 中没有最近点。
说明\cref{b1:thm:C1-distance-dual}的最大值在此仍可取得。
\end{enumerate}
\end{exercise}
% 来源：自拟；调用\chapref{b1:ch:09}已有范数不取得反例，新增最近点失败与对偶取得性并存的对比。
\begin{solution}
第一部分中，
\[
 \|(0,1)-(t,0)\|_\infty=\max\{|t|,1\}.
\]
所以距离为 \(1\)，全部最近点为 \(\{\,(t,0):|t|\leq1\,\}\)，并不唯一。

第二部分中，\cref{b1:ex:09-norm-attainment}已经证明 \(\|f\|=1\)，且
\(f\) 不在 \(c_0\) 的闭单位球上取得范数。因为 \(f(x)=1\)，有
\(\|x-c\|\geq f(x-c)=1\) 对每个 \(c\in C\) 成立。
令 \(v_N=\sum_{j=1}^Ne_j\)，则 \(f(v_N)=1-2^{-N}\)。
向量 \(c_N=x-v_N/(1-2^{-N})\) 属于 \(C\)，且
\[
 \|x-c_N\|_\infty=(1-2^{-N})^{-1}\longrightarrow1.
\]
故距离等于 \(1\)。若 \(c\) 为最近点，则 \(z=x-c\) 满足
\(\|z\|=1\)、\(f(z)=1\)，违反范数不取得，故不存在最近点。
然而对偶距离公式直接用单位泛函 \(f\) 就得到
\(f(x)-\sup_Cf=1\)，其最大值确实取得。
\end{solution}

\begin{exercise}\label{b1:ex:C1-finite-inequality-alternative}
给定实矩阵 \(A\in\mathbb R^{m\times n}\) 与 \(b\in\mathbb R^m\)。
证明以下条件等价：
\begin{equivalences}
\item\label{b1:item:C1-positive-solution}存在 \(x\in\mathbb R^n\)，使 \(Ax=b\) 且每个 \(x_j\geq0\)；
\item\label{b1:item:C1-positive-test}对每个满足 \(A^{\mathsf T}z\geq0\) 的 \(z\in\mathbb R^m\)，均有 \(z\cdot b\geq0\)。
\end{equivalences}
不预先引用多面体理论；须证明由 \(A\) 的列向量的非负组合组成的集合是闭集。
\end{exercise}
% 来源：自拟；把凸分离变成有限线性约束的对偶判别，关键增量是有限生成锥的闭性而非一句调用分离。
\begin{solution}
将 \(A\) 的列写成 \(a_1,\ldots,a_n\)，并令
\[
 C=\left\{\,\sum_{j=1}^nt_ja_j:t_j\geq0\,\right\}.
\]
先说明每个 \(y\in C\) 都有仅用线性无关列向量的表示。
若一个表示的正系数所对应的列线性相关，取这些列的一条非零关系
\(\sum_jc_ja_j=0\)，必要时整体变号，使某个 \(c_j>0\)。
令
\[
 s=\min_{c_j>0}\frac{t_j}{c_j}>0.
\]
改用 \(t_j-sc_j\) 后，全部系数仍非负，和不变，且至少一个正系数消失。
有限次重复即得所需表示。

现设 \(y_k\in C\)、\(y_k\to y\)。独立列的指标子集只有有限多个，可取子列使它们
使用同一个子集 \(I\)。若子集为空，子列恒为零，结论成立。
否则从 \(\mathbb R^{I}\) 到这些列张成空间的线性映射为有限维同构，
故其逆连续。相应系数向量随 \(y_k\) 收敛，极限各坐标非负，给出 \(y\in C\)。
因此 \(C\) 闭。

若 \(Ax=b\)、\(x\geq0\)，则
\(z\cdot b=(A^{\mathsf T}z)\cdot x\geq0\)，得到一个方向。
反过来，若 \(b\notin C\)，对闭凸集 \(C\) 与点 \(b\) 作严格分离，
可取线性泛函 \(h\)，使 \(\sup_Ch<h(b)\)。
由于 \(C\) 是含零且对非负数乘封闭的锥，这个有限上确界必为零，
并且 \(h(a_j)\leq0\) 对所有 \(j\) 成立。
写 \(h(v)=-z\cdot v\)，则 \(A^{\mathsf T}z\geq0\)，但 \(z\cdot b<0\)，
违反第二项。因此 \(b\in C\)，即存在所需的非负解。
\end{solution}
